
\chapter*{Preface}
\addcontentsline{toc}{chapter}{Preface}

A mathematical notebook is a very honest witness, though not always a reliable one. It preserves the order in which things became interesting, which is seldom the order in which they ought to be taught. A theorem may first appear as a name to be remembered, a definition as an obstacle to be crossed, and a proof as an inconvenience standing between the reader and the exercise. Years later the same page is opened again, and the theorem has become a method, the definition has acquired a necessity, and the proof has quietly taken possession of the subject. The notebook has not changed; the reader has.

The papers gathered here were written over a long interval and under many different provocations. Some began with a problem that refused to yield; some with a lecture, a book, a conversation, or a diagram that seemed to promise more than it explained. Others arose from the ordinary discipline of examinations, where one learns, among other useful things, how quickly an elegant idea can disappear when there are forty minutes left. None of these notes was composed with a future book in view. They accumulated by habit, and perhaps also by a certain distrust of understanding that leaves no written trace.

When such papers are brought together, one is tempted to arrange them into a history more orderly than the truth. The early calculation becomes the seed of a later theory; every confusion is promoted into a profound question; every detour is discovered to have been necessary. This is a flattering account, and for that reason should be treated with suspicion. Much of mathematical learning is accidental, repetitive, and badly timed. One meets an abstraction too early, forgets a useful lemma, proves by force what a better language would make immediate, and understands the example only after the general theory has become familiar. I have not tried to disguise these irregularities. They are part of the record, and often the most instructive part.

There is, however, a genuine continuity beneath them. An elementary calculation is frequently the first place where a larger structure shows itself, though at that moment it has neither a name nor a reputation. The Euclidean algorithm begins as a method of division and later speaks the language of ideals. A determinant is learned as a formula, then encountered again as volume, orientation, and the action of a linear map on an exterior power. A line integral, once understood through parametrized curves, becomes the pullback of a differential form. Fourier coefficients cease to be a table of numbers and become coordinates in an inner-product space. These later descriptions do not dignify the earlier work; the earlier work was already carrying them.

The traffic must also run in the opposite direction. Generality has a natural tendency to admire itself, and mathematics is not always improved by permitting it. A definition earns its place when it returns to an example and makes something clearer there. For this reason the present book keeps many computations, diagrams, special cases, and local arguments that a more economical account might dismiss. They are not merely illustrations placed after the theory has done its work. Very often they are where the theory first becomes visible, and where one can still tell whether an abstraction explains a phenomenon or has only renamed it.

\newpage

The same consideration applies to the circumstances in which a note was written. A seminar, a passing remark, or an image found at the right moment may determine the question more decisively than a formal syllabus. Such occasions deserve to remain when they belong to the movement of the mathematics. They should not, on the other hand, be preserved as relics merely because they happened. I have tried to keep the former and remove the latter: to let a question retain the situation that made it urgent, without requiring the reader to inspect the furniture of the room in which it occurred.

This work of revision has therefore been less a matter of polishing than of judgment. Errors had to be corrected, missing arguments supplied, and connections made explicit; but a perfectly smooth account would have been another kind of falsification. The mind does not proceed by numbered lemmas from ignorance to knowledge. It advances, retreats, changes notation, mistakes a resemblance for an identity, and occasionally reaches the right conclusion for reasons which do not survive inspection. Where a false start reveals something about the problem, I have allowed it to remain in a more disciplined form. Where it reveals only haste, it has been removed without ceremony.

The result is neither a textbook nor a diary. A textbook is expected to know in advance where it is going; a diary is under no such obligation. These notes stand somewhere between the two. They have been arranged so that they may be read, but not so completely arranged that they cease to resemble the way mathematics was encountered. Some chapters explain, some investigate, and some return to an old object with a language acquired elsewhere. The recurrence is deliberate. A subject seen twice is not always the same subject, and the second sight may owe as much to the first confusion as to the later theory.

Nor is the word \emph{informal} intended as an indulgence. An informal argument may still be exact, and a friendly tone does not excuse a careless statement. Definitions should bear the weight placed upon them; proofs should be checkable; examples should do more than decorate the page. What informality permits is a different relation between these things. A calculation may precede the theorem it suggests. A question may remain open at the end of a chapter. A difficult idea may be approached more than once, and from a side entrance. The object is not to pronounce the last word on any subject, but to leave enough of the path visible that the reader may continue it.

I have always thought that a useful note should survive the occasion that produced it without pretending to have escaped from time altogether. It should remember why a question mattered, yet remain intelligible when that occasion is forgotten. It should be capable of correction without losing its first energy, and of enlargement without becoming ashamed of its beginnings. The pages that follow are offered in that spirit. They are less a completed account of mathematics than a record of repeated approaches to it: sometimes direct, sometimes circuitous, but governed throughout by the belief that an idea understood once can still be understood better.

\begin{flushright}
Runnel Zhang\\
\today
\end{flushright}
